% Nouveau EGA V (cote 25): archive cote 25, site page 168 (PDF page 169).
% typescript. Draft compared: none.
% AI-transcribed by Claude Opus 5.5 (claude-opus-5-5, adaptive thinking, effort high), Message Batches API, batch ega25b (msgbatch_01A7n5YuXubUPNK1zsepi8RG), 2026-11-13; not yet checked by eye.
\scanpage{25/168}
\mspage{3}

EGA IV ..... Ensemble singulier et\xout{x} supersingulier.

Il me semble que le changement de point de vue possible signalé dans ma
dernière lettre peut-être récupéré après-coup, et qu'il peut y avoir intérêt
à partir d'abord de la situation où on a  un homomorphisme $\underline{E}\to\underline{O}_X$ , d'où
un $V^o = V(u) = V(u(\underline{E})\underline{O}_X)$, etc. Une fois cette étude faite, on fera les
remarques suivantes: supposons donné un X lisse(de dimension relative \hand{d} ..)
sur Y, un sous-préschéma $V^o = V(\underline{J})$ de X, tel que $\underline{J}$\xout{$/\underline{J}^2$xxxxNxxx} \xout{\unsure{K$\otimes$}xxxxxxxx} soit engendrab-%
le localement \inserted{\unsure{en les pts de $V^o$}} par m éléments (\unsure{$m \geq 0$} étant donné), où ce qui revient au
même, tel que $V^o$ soit localement de présentation finie sur X (ou encore sur
Y) et que \inserted{$\underline{N} =$} $\underline{N}_{V^o/X} = \underline{J}/\underline{J}^2$  soit localement engendrable par m éléments.  L'hy-%
pothèse peut aussi s'exprimer en disant que localement sur X, $\underline{J}$ peut être
représenté comme image d'un homomorphisme  u: $\underline{E}\to\underline{O}_X$, $\underline{E}$ étant localement
libre (ou même libre, tant qu'on \xout{m}reste dans le local !) de rang m.  En d'aut\unsure{x|r}
res termes, localement on est dans la situation de départ envisagée dans
les pages 14 et suivantes. Utilisant une telle représentation, on construit
alors des \inserted{sous-}préschémas \inserted{fermés} $V^1 \subset V^o$ , $V'' \subset V^1$, $V^2 \subset V^{\hand{1}}-V''$ . Il faut alors noter
(en proposition en forme) que ces sous-préschémas (et à fortiori les ensem-%
bles sous-jacents) sont en fait indépendants du choix de u:$\underline{E}\to\underline{O}_X$ définiss-%
ant $V^o$, \deleted{mais} \inserted{et} dépendent seu\over{\unsure{x}}{l}ement de \deleted{\unsure{X/Xxx}} $V^o$ ( et bien sur de X/Y et de
m !).  En effet, considérons d'abord  l'homomorphisme
\[ (\hand{$*^o$}) \qquad \underline{N}\otimes\underline{O}_{V^o}\to\underline{\Omega}^1_{X/Y}\otimes\underline{O}_{V^o} \ , \]
d'où  un homomorphisme
\[ (\hand{$*_m$}) \qquad \bigwedge^m\underline{N}\otimes\underline{O}_{V^o}\to\underline{\Omega}^{\over{\unsure{1}}{m}}_{X/Y}\otimes\underline{O}_{V^o}\ ; \]
on constate aussitôt que $V^1$ est le préschéma des zéros de cet homomorphisme
(qui ``existe'' grâce au fait que $\underline{\Omega}^{\hand{m}}_{X/Y}\otimes\underline{O}_{\unsure{V^o}}$  sur  $V^o$ est localement libre
de rang fini, et se définit comme réprésentant le foncteur qu'on devine
(Sch)$^o_{\hand{\unsure{/V^o}}}$ $\longrightarrow$ (Ens) ; c'est aussi le sous-préschéma de \over{V}{\unsure{$V^o$}}

\margin{168}
